\chapter{Introducción, Motivación y Objetivos}
\label{cap:introduccion}

Este trabajo estudia la viabilidad de incorporar representaciones cuánticas en el problema de construcción de carteras, manteniendo un enfoque reproducible y comparable frente a métodos clásicos consolidados.

\section{Motivación y objetivos}
\label{sec:motivacion_objetivos}

El problema de optimización de portafolios ha sido ampliamente estudiado en la literatura financiera. Sin embargo, con la llegada de la computación cuántica, surgen nuevas oportunidades para abordar estos problemas desde una perspectiva diferente.


La computación cuántica emerge como una alternativa prometedora para resolver problemas de optimización complejos. Algoritmos como el QAOA (Quantum Approximate Optimization Algorithm) y VQE (Variational Quantum Eigensolver) ofrecen potencial para acelerar cálculos que serían intratables en computadores clásicos. En particular, la capacidad de los circuitos cuánticos para codificar datos en espacios de alta dimensión mediante mapeos de características cuánticas permite capturar relaciones no lineales entre activos financieros que las métricas clásicas, basadas en la correlación lineal, no logran reflejar.

No obstante, aunque QAOA y VQE son relevantes en el panorama de optimización cuántica, este trabajo sigue un enfoque distinto basado en la codificación explícita de características y un ansatz variacional concreto para QHRP.


\section{Objetivos}
\label{sec:objetivos}

Los objetivos principales de este trabajo son:

\begin{enumerate}
\item Analizar la viabilidad de la codificación angular y del ansatz variacional en la representación de correlaciones financieras.
\item Implementar y reproducir el algoritmo Quantum Hierarchical Risk Parity (QHRP) propuesto por Arian et al.\ (2025), comparando los resultados con el enfoque clásico HRP.
\item Analizar el impacto de la representación cuántica sobre la estructura jerárquica del mercado y sobre el rendimiento de las carteras resultantes.
\end{enumerate}


\section{Estructura del documento}
\label{sec:estructura_documento}

El presente documento se organiza de la siguiente manera:

\begin{itemize}
\item El \textbf{Capítulo~\ref{cap:estado_arte}} reúne el contexto técnico y bibliográfico del trabajo, incluyendo fundamentos de computación cuántica y la evolución de los métodos clásicos y cuánticos en optimización de carteras.
\item El \textbf{Capítulo~\ref{cap:metodologia}} presenta la metodología empleada, incluyendo la representación multivariante de activos, el mapeo cuántico y la métrica de distancia utilizada.
\item El \textbf{Capítulo~\ref{cap:analisis_diseno}} desarrolla el análisis y diseño del pipeline, justificando las decisiones de modelado y arquitectura.
\item El \textbf{Capítulo~\ref{cap:desarrollo}} describe la implementación práctica, la reproducción de circuitos y la trazabilidad experimental del sistema.
\item El \textbf{Capítulo~\ref{cap:experimentos_resultados}} presenta los resultados experimentales y su análisis comparativo.
\item El \textbf{Capítulo~\ref{cap:conclusiones}} recoge las conclusiones del trabajo y las líneas de trabajo futuro.
\end{itemize}
